% !TEX program = xelatex
% The research corpus (ci/research.json): a Chinese paper in ctexart
% (xeCJK, TeX Live's Fandol fonts), with a Japanese line in IPAex.
\documentclass[UTF8,fontset=fandol,zihao=-4]{ctexart}
\usepackage{amsmath}
\usepackage{booktabs}
\setCJKfamilyfont{jp}{IPAexMincho}
\newcommand{\jp}{\CJKfamily{jp}}
\title{并行排版系统的设计与实现}
\author{张三 \and 李四}
\date{2026年1月1日}
\begin{document}
\maketitle
\begin{abstract}
本文介绍一种并行排版系统。文档被表示为一个依赖图，图中的节点可以并行执行，
其输出与串行排版逐字节相同。
\end{abstract}
\section{引言}
排版是一个古老的问题。\TeX{} 由高德纳（Donald E. Knuth）设计，至今仍被广泛使用。
中文排版需要处理标点挤压、中西文间距等问题，例如 PhiTeX 与 \LaTeX{} 混排。
\section{方法}
设文档有 $n$ 个段落，串行排版的时间为 $T_1 = \sum_{i=1}^{n} t_i$，
并行排版的时间为
\begin{equation}
  T_p \ge \max\Bigl(\frac{T_1}{p}, T_\infty\Bigr),
\end{equation}
其中 $T_\infty$ 为关键路径的长度。
\section{实验}
\begin{table}[h]
\centering
\caption{排版时间（秒）}
\begin{tabular}{lrr}
\toprule
文档 & 页数 & 时间 \\
\midrule
论文 & 12 & 0.8 \\
学位论文 & 180 & 6.4 \\
\bottomrule
\end{tabular}
\end{table}
\section{日本語}
{\jp 日本語の文章も同じ文書に含めることができる。「かぎ括弧」と句読点。}
\section{结论}
并行排版是可行的；其结果与串行排版完全一致。
\end{document}
